\begin{exercises}
  \item
    写出下列各向量空间中的零向量。
    \begin{exparts}
      \partsitem 在通常运算下，三次多项式构成的空间。
      \partsitem \( \nbym{2}{4} \) 矩阵构成的空间。
      \partsitem 空间
        \( \set{\map{f}{[0..1]}{\Re}\suchthat f\text{ 连续}} \)。
      \partsitem 一个自然数自变量的实值函数构成的空间。
    \end{exparts}
    \begin{answer}
      \begin{exparts}
        \partsitem \( 0+0x+0x^2+0x^3 \)
        \partsitem \( \begin{mat}[r]
                   0  &0  &0  &0  \\
                   0  &0  &0  &0
                 \end{mat} \)
        \partsitem 常值函数 \( f(x)=0 \)
        \partsitem 常值函数 \( f(n)=0 \)
      \end{exparts}
    \end{answer}
  \recommended \item
    求该向量在此向量空间中的加法逆元。
    \begin{exparts}
      \partsitem 在 \( \polyspace_3 \) 中，向量 \( -3-2x+x^2 \)。
      \partsitem 在空间 \( \nbyn{2} \) 中，
        \begin{equation*}
          \begin{mat}[r]
            1  &-1  \\
            0  &3
          \end{mat}.
        \end{equation*}
      \partsitem 在 \( \set{ae^x+be^{-x}\suchthat a,b\in\Re} \)——这是实变量
        \( x \) 的函数在通常运算下构成的空间——中，向量 \( 3e^x-2e^{-x} \)。
    \end{exparts}
    \begin{answer}
      \begin{exparts*}
        \partsitem \( 3+2x-x^2 \)
        \partsitem \( \begin{mat}[r]
                   -1  &+1  \\
                    0  &-3
                 \end{mat} \)
        \partsitem \( -3e^x+2e^{-x} \)
      \end{exparts*}
    \end{answer}
  \recommended \item
    对每一小题，列出三个元素，然后证明它是向量空间。
    \begin{exparts}
      \partsitem 一次多项式的集合
        \( \polyspace_1=\set{a_0+a_1x\suchthat a_0,a_1\in\Re} \)，
        运算为通常的多项式加法与标量乘法。
      \partsitem 一次多项式的集合
        \( \set{a_0+a_1x\suchthat a_0-2a_1=0} \)，
        运算为通常的多项式加法与标量乘法。
   \end{exparts}
   \textit{提示。}  以 \nearbyexample{PlaneThruOriginSubsp} 作为指导。
   十个条件中的大多数都只是验证而已。
   \begin{answer}
      \begin{exparts}
        \partsitem
          三个元素是：
          $1+2x$、$2-1x$ 与 $x$。
          （当然，可以有许多不同的答案。）

          验证过程与 \nearbyexample{ex:RealVecSpaces} 完全类似。
          我们先处理
          \nearbydefinition{def:VecSpace} 中与加法有关的
          条件 $1$-$5$。
          关于对加法封闭，即条件~(1)，
          注意当 $a+bx,c+dx\in\polyspace_1$ 时，
          $(a+bx)+(c+dx)=(a+c)+(b+d)x$ 是系数为实数的一次多项式，
          因而是
          $\polyspace_1$ 的元素。
          条件~(2) 的验证如下：当 $a+bx,c+dx\in\polyspace_1$ 时，
          $(a+bx)+(c+dx)=(a+c)+(b+d)x$，而按另一顺序相加则得
          $(c+dx)+(a+bx)=(c+a)+(d+b)x$，由于这些都是实数，
          故 $a+c=c+a$
          且 $b+d=d+b$。
          条件~(3) 类似：设
          $a+bx,c+dx,e+fx\in\polyspace$，则
          $((a+bx)+(c+dx))+(e+fx)=(a+c+e)+(b+d+f)x$
          而
          $(a+bx)+((c+dx)+(e+fx))=(a+c+e)+(b+d+f)x$，两者相等
          （也就是说，实数加法满足结合律，故 $(a+c)+e=a+(c+e)$
          且 $(b+d)+f=b+(d+f)$）。
          对于条件~(4)，注意一次多项式
          $0+0x\in\polyspace_1$ 满足
          $(a+bx)+(0+0x)=a+bx$ 与
          $(0+0x)+(a+bx)=a+bx$。
          对于这一部分的最后一个条件，即条件~(5)，
          注意对任意 $a+bx\in\polyspace_1$，其加法逆元
          是 $-a-bx\in\polyspace_1$，因为
          $(a+bx)+(-a-bx)=(-a-bx)+(a+bx)=0+0x$。

          接下来我们还要检查与
          标量乘法有关的条件~(6)-(10)。
          对于~(6)，即空间对标量乘法封闭这一条件，
          设 $r$ 是实数，$a+bx$ 是
          $\polyspace_1$ 的元素，则 $r(a+bx)=(ra)+(rb)x$ 是
          $\polyspace_1$ 的元素，因为它是系数为
          实数的一次多项式。
          条件~(7) 成立，因为
          $(r+s)(a+bx)=r(a+bx)+s(a+bx)$ 可由实数乘法的分配律
          得到。
          条件~(8) 类似：
          $r((a+bx)+(c+dx))=r((a+c)+(b+d)x)=r(a+c)+r(b+d)x=(ra+rc)+(rb+rd)x
          =r(a+bx)+r(c+dx)$。
          对于~(9)，我们有
          $(rs)(a+bx)=(rsa)+(rsb)x=r(sa+sbx)=r(s(a+bx))$。
          最后，条件~(10) 是
          $1(a+bx)=(1a)+(1b)x=a+bx$。
        \partsitem
          将该集合记为 $P$。
          本练习的前一小题对系数没有限制，
          而这里我们只关注满足 $a_0-2a_1=0$ 的一次多项式，
          即常数项减去一次项系数的两倍为零的一次多项式。
          于是，$P$ 的三个典型元素是
          $2+1x$、$6+3x$ 与 $-4-2x$。

          对于条件~(1)，我们必须证明：
          若两个满足该限制的一次多项式相加，
          则所得仍是一次满足该限制的一次多项式：这里的论证是，
          若 $a+bx,c+dx\in P$，则
          $(a+bx)+(c+dx)=(a+c)+(b+d)x$ 是 $P$ 的元素，因为
          $(a+c)-2(b+d)=(a-2b)+(c-2d)=0+0=0$。
          条件~(2) 可以这样验证：
          当 $a+bx,c+dx\in\polyspace_1$ 时，
          $(a+bx)+(c+dx)=(a+c)+(b+d)x$，而按另一顺序相加则得
          $(c+dx)+(a+bx)=(c+a)+(d+b)x$，由于这些都是实数，
          故 $a+c=c+a$
          且 $b+d=d+b$。
          （也就是说，这一条件不受该限制的影响，
          其验证与本练习第一小题中的验证相同。）
          条件~(3) 同样不受这一附加限制的影响：
          设
          $a+bx,c+dx,e+fx\in\polyspace$，则
          $((a+bx)+(c+dx))+(e+fx)=(a+c+e)+(b+d+f)x$
          而
          $(a+bx)+((c+dx)+(e+fx))=(a+c+e)+(b+d+f)x$，两者相等。
          对于条件~(4)，注意一次多项式
          $0+0x\in P$ 满足该限制，因为
          其常数项减去其一次项系数的两倍
          为零，于是本小题中的
          验证即可适用：
          $0+0x\in\polyspace_1$ 满足
          $(a+bx)+(0+0x)=a+bx$ 与
          $(0+0x)+(a+bx)=a+bx$。
          为检查条件~(5)，
          注意对任意 $a+bx\in P$，其加法逆元
          是 $-a-bx$，因为
          它是 $P$ 的元素（由 $a+bx\in P$ 知
          $a-2b=0$，两边同乘 $-1$ 得 $-a+2b=0$），
          并且如第一小题那样，它起到加法逆元的作用：
          $(a+bx)+(-a-bx)=(-a-bx)+(a+bx)=0+0x$。

          我们还必须检查
          标量乘法的条件~(6)-(10)。
          对于~(6)，即空间对标量乘法封闭这一条件，
          设 $r$ 是实数且 $a+bx\in P$（于是 $a-2b=0$），
          则 $r(a+bx)=(ra)+(rb)x$ 是
          $P$ 的元素，因为它是系数为实数、
          且满足 $(ra)-2(rb)=r(a-2b)=0$ 的一次多项式。
          条件~(7) 成立的理由与本练习第一小题中
          相同，因为
          $(r+s)(a+bx)=r(a+bx)+s(a+bx)$ 可由实数乘法的分配律
          得到。
          条件~(8) 也与第一小题中相同：
          $r((a+bx)+(c+dx))=r((a+c)+(b+d)x)=r(a+c)+r(b+d)x=(ra+rc)+(rb+rd)x
          =r(a+bx)+r(c+dx)$。
          对~(9) 也是如此：
          $(rs)(a+bx)=(rsa)+(rsb)x=r(sa+sbx)=r(s(a+bx))$。
          最后，条件~(10) 也是如此：
          $1(a+bx)=(1a)+(1b)x=a+bx$。
       \end{exparts}
     \end{answer}
  \item
    对每一小题，列出三个元素，然后证明它是向量空间。
    \begin{exparts}
      \partsitem 实元素的 \( \nbyn{2} \) 矩阵的集合，运算为通常的矩阵运算。
      \partsitem 实元素且 $2,1$ 元为零的 \( \nbyn{2} \) 矩阵的集合，
        运算为通常的矩阵运算。
    \end{exparts}
    \begin{answer}
      可参照
      \nearbyexample{ex:RealVecSpaces} 作为指导。
      （\textit{评注。}
      由于许多条件都很容易检查，
      有时人们会觉得一定是漏掉了什么。
      请记住：容易做或常规，与不必做是两回事。）
      \begin{exparts}
        \partsitem
          以下是三个元素。
          \begin{equation*}
            \begin{mat}
              1  &2  \\
              3  &4
            \end{mat},\,
            \begin{mat}
              -1  &-2  \\
              -3  &-4
            \end{mat},\,
            \begin{mat}
              0  &0  \\
              0  &0
            \end{mat}
          \end{equation*}

          对于~(1)，$\nbyn{2}$ 实矩阵的和是 $\nbyn{2}$
          实矩阵。
          对于~(2)，我们考察两个矩阵的和
          \begin{equation*}
            \begin{mat}
              a  &b  \\
              c  &d
            \end{mat}+
            \begin{mat}
              e  &f  \\
              g  &h
            \end{mat}
            =
            \begin{mat}
              a+e  &b+f  \\
              c+g  &d+h
            \end{mat}
          \end{equation*}
          并利用实数加法的交换律
          \begin{equation*}
            =\begin{mat}
              e+a  &f+b  \\
              g+c  &h+d
            \end{mat}
            =
            \begin{mat}
              e  &f  \\
              g  &h
            \end{mat}
            +
            \begin{mat}
              a  &b  \\
              c  &d
            \end{mat}
          \end{equation*}
          来验证矩阵加法的交换性。
          条件~(3) 的验证，即矩阵加法的结合律，
          与前面的验证类似：
          \begin{equation*}
            \big(
              \begin{mat}
                a  &b  \\
                c  &d
              \end{mat}
              +\begin{mat}
                e  &f  \\
                g  &h
              \end{mat}
            \big)
            +
            \begin{mat}
              i  &j  \\
              k  &l
            \end{mat}
            =
            \begin{mat}
              (a+e)+i  &(b+f)+j  \\
              (c+g)+k  &(d+h)+l
            \end{mat}
          \end{equation*}
          而
          \begin{equation*}
            \begin{mat}
              a  &b  \\
              c  &d
            \end{mat}
            +
              \big(
              \begin{mat}
                e  &f  \\
                g  &h
              \end{mat}
              +
              \begin{mat}
                i  &j  \\
                k  &l
              \end{mat}
            \big)
            =
            \begin{mat}
              a+(e+i)  &b+(f+j)  \\
              c+(g+k)  &d+(h+l)
            \end{mat}
          \end{equation*}
          且由于实数加法满足结合律，两者逐项相同。
          对于~(4)，这个空间的零元素是全为零的 $\nbyn{2}$
          矩阵。
          条件~(5) 成立，因为对任意 $\nbyn{2}$ 矩阵~$A$，
          其加法逆元是各元素为 $A$ 相应元素的相反数的矩阵，
          即矩阵 $-1\cdot A$。

          条件~($6$)
          成立，因为 $\nbyn{2}$ 矩阵的标量倍
          仍是 $\nbyn{2}$ 矩阵。
          对于条件~(7)，我们有：
          \begin{multline*}
            (r+s)
            \begin{mat}
              a  &b  \\
              c  &d
            \end{mat}
            =
            \begin{mat}
              (r+s)a  &(r+s)b  \\
              (r+s)c  &(r+s)d
            \end{mat}               \\
            =
            \begin{mat}
              ra+sa  &rb+sb  \\
              rc+sc  &rd+sd
            \end{mat}
            =
            r
            \begin{mat}
              a  &b  \\
              c  &d
            \end{mat}
            +
            s
            \begin{mat}
              a  &b  \\
              c  &d
            \end{mat}
          \end{multline*}
          条件~(8) 的处理方式相同。
          \begin{multline*}
            r
            \big(
              \begin{mat}
                a  &b  \\
                c  &d
              \end{mat}
              +
              \begin{mat}
                e  &f  \\
                g  &h
              \end{mat}
            \big)
            =
            r
            \begin{mat}
              a+e  &b+f  \\
              c+g  &d+h
            \end{mat}
            =
            \begin{mat}
              ra+re  &rb+rf  \\
              rc+rg  &rd+rh
            \end{mat}              \\
            =
            r
            \begin{mat}
              a  &b  \\
              c  &d
            \end{mat}
            +
            r
            \begin{mat}
              e  &f  \\
              g  &h
            \end{mat}
            =
            r
            \big(
            \begin{mat}
              a  &b  \\
              c  &d
            \end{mat}
            +
            \begin{mat}
              e  &f  \\
              g  &h
            \end{mat}
            \big)
          \end{multline*}
          对于~(9)，我们有：
          \begin{equation*}
            (rs)
            \begin{mat}
              a  &b  \\
              c  &d
            \end{mat}
            =
            \begin{mat}
              rsa  &rsb  \\
              rsc  &rsd
            \end{mat}
            =
            r
            \begin{mat}
              sa  &sb  \\
              sc  &sd
            \end{mat}
            =
            r\big( s
            \begin{mat}
              a  &b  \\
              c  &d
            \end{mat}
            \big)
          \end{equation*}
          条件~(10) 同样简单。
          \begin{equation*}
            1
            \begin{mat}
              a  &b  \\
              c  &d
            \end{mat}
            =
            \begin{mat}
              1\cdot a  &1\cdot b  \\
              1\cdot c  &1\cdot d
            \end{mat}
            =
            \begin{mat}
              sa  &sb  \\
              sc  &sd
            \end{mat}
          \end{equation*}
        \partsitem
          本小题与本练习前一小题的差别仅在于：
          我们限定于第二行第一列
          为零的矩阵集合 $T$。
          以下是 $T$ 的三个元素。
          \begin{equation*}
            \begin{mat}
              1  &2  \\
              0  &4
            \end{mat},\,
            \begin{mat}
              -1  &-2  \\
              0  &-4
            \end{mat},\,
            \begin{mat}
              0  &0  \\
              0  &0
            \end{mat}
          \end{equation*}
          本小题的某些验证与本练习第一小题相同，
          下面我们只做那些不同的验证。

          对于~(1)，$2,1$ 元为零的 $\nbyn{2}$ 实矩阵之和
          也是 $2,1$ 元为零的 $\nbyn{2}$ 实矩阵。
          \begin{equation*}
            \begin{mat}
              a  &b  \\
              0  &d
            \end{mat}
            +
            \begin{mat}
              e  &f  \\
              0  &h
            \end{mat}
            \begin{mat}
              a+e  &b+f  \\
              0    &d+h
            \end{mat}
          \end{equation*}
          前一小题中给出的条件~(2) 的验证
          在本小题同样适用。
          条件~(3) 也是如此。
          对于~(4)，注意全为零的 $\nbyn{2}$
          矩阵是 $T$ 的元素。
          条件~(5) 成立，因为对任意 $\nbyn{2}$ 矩阵~$A$，
          其加法逆元是矩阵
          $-1\cdot A$，因此 $2,1$ 元为零的矩阵的加法逆元
          也是 $2,1$ 元为零的矩阵。

          条件~$6$ 成立，因为 $2,1$ 元为零的 $\nbyn{2}$ 矩阵的标量倍
          仍是 $2,1$ 元为零的 $\nbyn{2}$ 矩阵。
          条件~(7) 的验证与前一小题相同。
          条件~(8)、(9) 与~(10) 的验证也是如此。
      \end{exparts}
    \end{answer}
  \recommended \item
    对每一小题，列出三个元素，然后证明它是向量空间。
    \begin{exparts}
      \partsitem 三个分量的行向量的集合，运算为其通常运算。
      \partsitem 集合
        \begin{equation*}
          \set{\colvec{x \\ y \\ z \\ w}\in\Re^4\suchthat x+y-z+w=0}
        \end{equation*}
        运算取从 $\Re^4$ 继承的运算。
    \end{exparts}
    \begin{answer}
      % Most of the conditions are easy to check; use
      % \nearbyexample{ex:RealVecSpaces} as a guide.
      \begin{exparts}
        \partsitem
          三个元素是 $\rowvec{1  &2  &3}$、
          $\rowvec{2  &1  &3}$ 与 $\rowvec{0  &0  &0}$。

          我们必须检查 \nearbydefinition{def:VecSpace} 中的条件 (1)-(10)。
          条件 (1)-(5) 与加法有关。
          对于条件~(1)，回顾两个三分量行向量的和
          \begin{equation*}
            \rowvec{a  &b  &c}
            +\rowvec{d  &e  &f}
            =\rowvec{a+d  &b+e  &c+f}
          \end{equation*}
          也是三分量行向量
         （字母 $a,\ldots,f$ 全都表示实数）。
         对~(2) 的验证是常规的
          \begin{multline*}
            \rowvec{a  &b  &c}
            +\rowvec{d  &e  &f}
            =\rowvec{a+d  &b+e  &c+f}   \\
            =\rowvec{d+a  &e+b  &f+c}
            =\rowvec{d  &e  &f}
            +\rowvec{a  &b  &c}
          \end{multline*}
          （第二个等号成立是因为三个分量都是实数，
          而实数加法可交换）。
         条件~(3) 的验证类似。
          \begin{multline*}
            \big(\rowvec{a  &b  &c}
              +\rowvec{d  &e  &f}
            \big)
            +\rowvec{g  &h  &i}
            =\rowvec{(a+d)+g  &(b+e)+h  &(c+f)+i}    \\
            =\rowvec{a+(d+g)  &b+(e+h)  &c+(f+i)}
            =\rowvec{a  &b  &c}
            +\big(\rowvec{d  &e  &f}
               +\rowvec{g  &h  &i}
             \big)
          \end{multline*}
          对于~(4)，注意三分量行向量
          $\rowvec{0  &0  &0}$ 是加法单位元：
          $\rowvec{a  &b  &c}+\rowvec{0  &0  &0}=\rowvec{a  &b  &c}$。
          为验证条件~(5)，设给定集合中的
          元素 $\rowvec{a  &b  &c}$，并注意
          $\rowvec{-a  &-b  &-c}$ 也在此集合中且具有
          所需性质：
          $\rowvec{a  &b  &c}+\rowvec{-a  &-b  &-c}=\rowvec{0  &0  &0}$。

          条件 (6)-(10) 涉及标量乘法。
          为验证~(6)，即空间对所给的
          标量乘法运算封闭，
          注意 $r\rowvec{a  &b  &c}=\rowvec{ra  &rb  &rc}$ 是
          实分量的三分量行向量。
          对于~(7)，我们计算如下。
          \begin{multline*}
            (r+s)\rowvec{a  &b  &c}=\rowvec{(r+s)a  &(r+s)b  &(r+s)c}
            =\rowvec{ra+sa  &rb+sb  &rc+sc}                             \\
            =\rowvec{ra  &rb  &rc}+\rowvec{sa  &sb  &sc}
            =r\rowvec{a  &b  &c}+s\rowvec{a  &b  &c}
          \end{multline*}
          条件~(8) 非常类似。
          \begin{multline*}
            r\big(\rowvec{a  &b  &c}+\rowvec{d  &e  &f}\big)  \\
            \begin{aligned}
            &=r\rowvec{a+d  &b+e  &c+f}
            =\rowvec{r(a+d)  &r(b+e)  &r(c+f)}                     \\
            &=\rowvec{ra+rd  &rb+re  &rc+rf}
            =\rowvec{ra  &rb  &rc}+\rowvec{rd  &re  &rf}          \\
            &=r\rowvec{a  &b  &c}+r\rowvec{d  &e  &f}
            \end{aligned}
          \end{multline*}
          条件~(9) 的计算也是如此。
          \begin{equation*}
            (rs)\rowvec{a  &b  &c}
             =\rowvec{rsa  &rsb  &rsc}
             =r\rowvec{sa  &sb  &sc}
             =r\big(s\rowvec{a  &b  &c}\big)
          \end{equation*}
          条件~(10) 同样是常规的：
          $1\rowvec{a  &b  &c}
             =\rowvec{1\cdot a  &1\cdot b  &1\cdot c}
             =\rowvec{a  &b  &c}$。
        \partsitem
          将该集合记为 $L$。
          加法的封闭性，即条件~(1)，需要检查：若各加数都是~$L$ 的成员，则
          和
          \begin{equation*}
            \colvec{a \\ b \\ c \\ d}
            +\colvec{e \\ f \\ g \\ h}
            =
            \colvec{a+e \\ b+f \\ c+g \\ d+h}
          \end{equation*}
          也是 \( L \) 的成员；这是成立的，因为它满足
          $L$ 的成员资格判据：
          \( (a+e)+(b+f)-(c+g)+(d+h)
          =(a+b-c+d)+(e+f-g+h)=0+0 \)。
          条件~(2)、(3) 与~(5) 的验证与本练习第一部分中的
          验证类似。
          对于条件~(4)，注意零向量是~$L$ 的成员，
          因为其第一分量加第二分量、减第三分量、
          再加第四分量，总和为零。

          条件~(6)，即标量乘法的封闭性，类似：
          当该向量是~$L$ 的元素时，
          \begin{equation*}
            r\colvec{a  \\ b \\ c \\ d}
            =\colvec{ra  \\  rb  \\  rc  \\  rd}
          \end{equation*}
          也是~$L$ 的元素，因为 $ra+rb-rc+rd=r(a+b-c+d)=r\cdot 0=0$。
          条件~(7)、(8)、(9) 与~(10) 的验证与本练习
          前一小题中相同。
     \end{exparts}
     \end{answer}
  \recommended \item \label{exer:NotVectorSpaces}
    证明下列每一个都不是向量空间。
    （\textit{提示。}  为检查封闭性，列出每个集合的两个成员，
    并对它们尝试一些运算。）
    \begin{exparts}
      \partsitem 在从 \( \Re^3 \) 继承的运算下，集合
        \begin{equation*}
          \set{\colvec{x \\ y \\ z}\in\Re^3\suchthat x+y+z=1}
        \end{equation*}
      \partsitem 在从 \( \Re^3 \) 继承的运算下，集合
        \begin{equation*}
          \set{\colvec{x \\ y \\ z}\in\Re^3\suchthat x^2+y^2+z^2=1}
        \end{equation*}
      \partsitem 在通常的矩阵运算下，
        \begin{equation*}
          \set{\begin{mat}
                 a  &1  \\
                 b  &c
               \end{mat} \suchthat a,b,c\in\Re}
        \end{equation*}
      \partsitem 在通常的多项式运算下，
        \begin{equation*}
          \set{a_0+a_1x+a_2x^2\suchthat a_0,a_1,a_2\in\Re^+}
        \end{equation*}
        其中 $\Re^+$ 是大于零的实数构成的集合
      \partsitem 在继承的运算下，
        \begin{equation*}
          \set{\colvec{x \\ y}\in\Re^2\suchthat
               \text{\( x+3y=4 \) 且 \( 2x-y=3 \) 且 \( 6x+4y=10 \)} }
        \end{equation*}
    \end{exparts}
    \begin{answer}
      每一小题中都将该集合记为 \( Q \)。
      对其中某些小题，还有其他正确方法可证明 $Q$ 不是
      向量空间。
      \begin{exparts}
        \partsitem 它对加法不封闭；它不满足
          条件~(1)。
          \begin{equation*}
            \colvec[r]{1 \\ 0 \\ 0},
            \colvec[r]{0 \\ 1 \\ 0}\in Q
            \qquad
            \colvec[r]{1 \\ 1 \\ 0}\not\in Q
          \end{equation*}
        \partsitem 它对加法不封闭。
          \begin{equation*}
            \colvec[r]{1 \\ 0 \\ 0},
            \colvec[r]{0 \\ 1 \\ 0}\in Q
            \qquad
            \colvec[r]{1 \\ 1 \\ 0}\not\in Q
          \end{equation*}
        \partsitem 它对加法不封闭。
          \begin{equation*}
            \begin{mat}[r]
              0  &1  \\
              0  &0
            \end{mat},
            \,
            \begin{mat}[r]
              1  &1  \\
              0  &0
            \end{mat}\in Q
            \qquad
            \begin{mat}[r]
              1  &2  \\
              0  &0
            \end{mat}\not\in Q
          \end{equation*}
        \partsitem 它对标量乘法不封闭。
          \begin{equation*}
            1+1x+1x^2\in Q
            \qquad
            -1\cdot(1+1x+1x^2)\not\in Q
          \end{equation*}
        \item 它是空集，违反条件~(4)。
      \end{exparts}
    \end{answer}
  \item
    定义加法与标量乘法运算，
    使复数成为 \( \Re \) 上的向量空间。
    \begin{answer}
      通常的运算
      \( (v_0+v_1i)+(w_0+w_1i)=(v_0+w_0)+(v_1+w_1)i \) 与
      \( r(v_0+v_1i)=(rv_0)+(rv_1)i \) 即可满足要求。
      检查很容易。
    \end{answer}
  \item
    在通常的加法与标量乘法
    运算下，有理数集是 \( \Re \) 上的向量空间吗？
    \begin{answer}
       不是，它对标量乘法不封闭，因为例如
       \( \pi\cdot (1) \) 不是有理数。
    \end{answer}
  \item
    证明
    变量 \( x,y,z \) 的线性组合构成的集合
    在自然的加法与标量乘法
    运算下是向量空间。
    \begin{answer}
      自然运算是
      \( (v_1x+v_2y+v_3z)+(w_1x+w_2y+w_3z)=(v_1+w_1)x+(v_2+w_2)y+(v_3+w_3)z \)
      与 \( r\cdot(v_1x+v_2y+v_3z)=(rv_1)x+(rv_2)y+(rv_3)z \)。
      检查它是向量空间很容易；以
      \nearbyexample{ex:RealVecSpaces} 作为指导。
    \end{answer}
  \item
    证明
    这不是向量空间：实元素的
    二行列向量的集合，其运算如下。
    \begin{equation*}
      \colvec{x_1 \\ y_1}
      +\colvec{x_2 \\ y_2}
      =\colvec{x_1-x_2 \\ y_1-y_2}
      \qquad
      r\cdot\colvec{x \\ y}
      =\colvec{rx \\ ry}
    \end{equation*}
    \begin{answer}
      “\( + \)” 运算不可交换（即条件~(2) 不
      满足）；很容易找出该集合的两个成员来
      证实这一断言。
    \end{answer}
  \item
    证明或反驳：\( \Re^3 \) 在下列
    运算下是向量空间。
    \begin{exparts}
      \partsitem \(
               \colvec{x_1 \\ y_1 \\ z_1}
               +\colvec{x_2 \\ y_2 \\ z_2}
               =\colvec{0 \\ 0 \\ 0}
               \quad\text{且}\quad
               r\colvec{x \\ y \\ z}
               =\colvec{rx \\ ry \\ rz} \)
      \partsitem \(
               \colvec{x_1 \\ y_1 \\ z_1}
               +\colvec{x_2 \\ y_2 \\ z_2}
               =\colvec{0 \\ 0 \\ 0}
               \quad\text{且}\quad
               r\colvec{x \\ y \\ z}
               =\colvec{0 \\ 0 \\ 0} \)
    \end{exparts}
    \begin{answer}
      \begin{exparts}
        \partsitem 它不是向量空间。
          \begin{equation*}
            (1+1)\cdot\colvec[r]{1 \\ 0 \\ 0}\neq
            \colvec[r]{1 \\ 0 \\ 0}
            +\colvec[r]{1 \\ 0 \\ 0}
          \end{equation*}
        \partsitem 它不是向量空间。
          \begin{equation*}
            1\cdot\colvec[r]{1 \\ 0 \\ 0}\neq\colvec[r]{1 \\ 0 \\ 0}
          \end{equation*}
      \end{exparts}
    \end{answer}
  \recommended \item
    对每一小题，判断它是否为向量空间；
    所指的运算是自然的那些。
    \begin{exparts}
      \partsitem \definend{对角} \( \nbyn{2} \) 矩阵
        \begin{equation*}
          \set{\begin{mat}
            a  &0  \\
            0  &b
          \end{mat}\suchthat a,b\in\Re}
        \end{equation*}
      \partsitem 这个 \( \nbyn{2} \) 矩阵集合
        \begin{equation*}
          \set{\begin{mat}
            x    &x+y  \\
            x+y  &y
          \end{mat}\suchthat x,y\in\Re}
        \end{equation*}
      \partsitem 这个集合
        \begin{equation*}
          \set{\colvec{x \\ y \\ z \\ w}\in\Re^4
               \suchthat x+y+z+w=1}
        \end{equation*}
      \partsitem 函数的集合
        \( \set{\map{f}{\Re}{\Re}\suchthat df/dx+2f=0} \)
      \partsitem 函数的集合
        \( \set{\map{f}{\Re}{\Re}\suchthat df/dx+2f=1} \)
    \end{exparts}
    \begin{answer}
      对每个“是”的答案，只需在该空间中检查几个
      线性组合即可。
      对每个“否”的答案，给出其中一个
      条件不成立的具体例子。
      \begin{exparts}
        \partsitem 是。
          随手取的一个线性组合如下。
          \begin{equation*}
            2\cdot\begin{mat}
              3  &0 \\
              0  &4
            \end{mat}
            +5\cdot\begin{mat}
              6  &0 \\
              0  &7
            \end{mat}
          \end{equation*}
          这个集合成为向量空间的关键在于：
          结果是对角矩阵。
          \begin{equation*}
            =\begin{mat}
              36 &0 \\
              0  &43
            \end{mat}
          \end{equation*}
        \partsitem 是。
          线性组合如下。
          \begin{equation*}
            10\cdot\begin{mat}
              1  &3  \\
              3  &2
            \end{mat}
            +11\cdot\begin{mat}
              3  &7  \\
              7  &4
            \end{mat}
          \end{equation*}
          我们看到结果保持了如下形式：左上
          与右下元素相加得到右上元素，
          同时也得到左下元素。
          \begin{equation*}
            =\begin{mat}
               43  &107  \\
              107  &64
            \end{mat}
          \end{equation*}
        \partsitem 否，这个集合对自然加法
          运算不封闭。
          全部分量都是 $1/4$ 的向量是这个集合的成员，
          但它自加所得的结果，即
          全部分量都是 $1/2$ 的向量，不是成员。
        \partsitem 是。
        \partsitem 否，\( f(x)=e^{-2x}+(1/2) \) 在集合中，
           但 \( 2\cdot f \) 不在（即条件~(6) 不成立）。
      \end{exparts}
    \end{answer}
  \recommended \item
    证明或反驳：这是一个向量空间——满足 \( f(7)=0 \) 的
    单实变量实值函数 \( f \) 构成的集合。
    \begin{answer}
      它是向量空间。
      向量空间定义中的大多数条件的验证都是常规的；这里我们
      只检查封闭性。
      对于加法，
      \( (f_1+f_2)\,(7)=f_1(7)+f_2(7)=0+0=0 \)。
      对于标量乘法，
      \( (r\cdot f)\,(7)=rf(7)=r0=0 \)。
    \end{answer}
  \recommended \item
    证明：正实数集 \( \Re^+ \)
    是向量空间，其中我们把 “\( x+y \)” 解释为
    \( x \) 与 \( y \) 的乘积（于是 \( 2+3 \) 为 \( 6 \)），
    把 “\( r\cdot x \)” 解释为 \( x \) 的 \( r \) 次幂。
    \begin{answer}
      我们逐条检查 \nearbydefinition{def:VecSpace}。

